Geothermics is unique among geophysical methods because it directly measures and models a thermodynamic state variable: temperature.
Whereas gravity, magnetics, electromagnetics, and seismology infer structure through proxy responses that depend on pressure- and temperature-dependent physical properties, geothermics observes temperature itself and models its evolution through time. This directness is also a limitation: thermal observations are integrative, diffusive, and do not image discrete structures in the same way as wave-based or potential-field methods.

Despite this lack of spatial imaging, temperature plays a central role in Earth systems. It controls the rates of deformation, diffusion, and chemical reactions; sets thresholds for melting, metamorphism, and rheological transitions; and governs the magnitude of properties such as viscosity, electrical conductivity, seismic attenuation, and elastic moduli. As a result, thermal structure exerts first-order control on tectonic behaviour, lithospheric strength, crustal evolution, and the longevity of geological provinces.

In geothermics, present-day temperature fields reflect both current boundary conditions and the integrated effects of past tectonic, magmatic, and erosional processes. Interpreting thermal data therefore requires careful attention to material properties — particularly thermal conductivity and radiogenic heat production — and to the distinction between steady-state and transient behaviour. More than any other method in this course, geothermics highlights the trade-offs between physical realism, mathematical simplicity, and interpretability.